% PDF-to-TeX transcription by the assistant; not the original downloaded file.

\begin{theorem}[4.54, Double Centralizer Theorem]
Let $A,B$ be two subalgebras of the algebra $\operatorname{End}E$ of endomorphisms of a finite dimensional vector space $E$, such that $A$ is semisimple, and $B=\operatorname{End}_A E$. Then:
\begin{enumerate}[label=(\roman*)]
\item $A=\operatorname{End}_B E$ (i.e., the centralizer of the centralizer of $A$ is $A$);
\item $B$ is semisimple;
\item as a representation of $A\otimes B$, $E$ decomposes as $E=\bigoplus_{i\in I}V_i\otimes W_i$, where $V_i$ are all the irreducible representations of $A$, and $W_i$ are all the irreducible representations of $B$. In particular, we have a natural bijection between irreducible representations of $A$ and $B$.
\end{enumerate}
\end{theorem}
\begin{proof}
Since $A$ is semisimple, we have a natural decomposition $E=\bigoplus_{i\in I}V_i\otimes W_i$, where $W_i:=\Hom_A(V_i,E)$, and $A=\bigoplus_i\operatorname{End}V_i$. Therefore, by Schur's lemma, $B=\operatorname{End}_A(E)$ is naturally identified with $\bigoplus_i\operatorname{End}(W_i)$. This implies all the statements of the theorem.
\end{proof}
We will now apply Theorem 4.54 to the following situation: $E=V^{\otimes n}$, where $V$ is a finite dimensional vector space over a field of characteristic zero, and $A$ is the image of $\C[S_n]$ in $\operatorname{End}E$. Let us now characterize the algebra $B$. Let $\mathfrak{gl}(V)$ be $\operatorname{End}V$ regarded as a Lie algebra with operation $ab-ba$.
\begin{theorem}[4.55]
The algebra $B=\operatorname{End}_A E$ is the image of the universal enveloping algebra $U(\mathfrak{gl}(V))$ under its natural action on $E$. In other words, $B$ is generated by elements of the form
\[
\Delta_n(b):=b\otimes1\otimes\cdots\otimes1+1\otimes b\otimes\cdots\otimes1+\cdots+1\otimes1\otimes\cdots\otimes b,
\quad b\in\mathfrak{gl}(V).
\]
\end{theorem}
\begin{proof}
Clearly, the image of $U(\mathfrak{gl}(V))$ is contained in $B$, so we just need to show that any element of $B$ is contained in the image of $U(\mathfrak{gl}(V))$. By definition, $B=S^n\operatorname{End}V$, so the result follows from part (ii) of the following lemma.
\begin{lemma}[4.56]
Let $k$ be a field of characteristic zero.
\begin{enumerate}[label=(\roman*)]
\item For any finite dimensional vector space $U$ over $k$, the space $S^nU$ is spanned by elements of the form $u\otimes\cdots\otimes u$, $u\in U$.
\item For any algebra $A$ over $k$, the algebra $S^nA$ is generated by elements $\Delta_n(a)$, $a\in A$.
\end{enumerate}
\end{lemma}
\begin{proof}
(i) The space $S^nU$ is an irreducible representation of $\operatorname{GL}(U)$ (Problem 3.19). The subspace spanned by $u\otimes\cdots\otimes u$ is a nonzero subrepresentation, so it must be everything.

(ii) By the fundamental theorem on symmetric functions, there exists a polynomial $P$ with rational coefficients such that $P(H_1(x),\ldots,H_n(x))=x_1\cdots x_n$ (where $x=(x_1,\ldots,x_n)$). Then
\[
P(\Delta_n(a),\Delta_n(a^2),\ldots,\Delta_n(a^n))=a\otimes\cdots\otimes a.
\]
The rest follows from (i).
\end{proof}
\end{proof}
Now, the algebra $A$ is semisimple by Maschke's theorem, so the double centralizer theorem applies, and we get the following result, which goes under the name ``Schur--Weyl duality''.
\begin{theorem}[4.57]
\begin{enumerate}[label=(\roman*)]
\item The image $A$ of $\C[S_n]$ and the image $B$ of $U(\mathfrak{gl}(V))$ in $\operatorname{End}(V^{\otimes n})$ are centralizers of each other.
\item Both $A$ and $B$ are semisimple. In particular, $V^{\otimes n}$ is a semisimple $\mathfrak{gl}(V)$-module.
\item We have a decomposition of $A\otimes B$-modules $V^{\otimes n}=\bigoplus_\lambda V_\lambda\otimes L_\lambda$, where the summation is taken over partitions of $n$, $V_\lambda$ are Specht modules for $S_n$, and $L_\lambda$ are some distinct irreducible representations of $\mathfrak{gl}(V)$ or zero.
\end{enumerate}
\end{theorem}
The Schur--Weyl duality for the Lie algebra $\mathfrak{gl}(V)$ implies a similar statement for the group $\operatorname{GL}(V)$.
\begin{proposition}[4.58]
The image of $\operatorname{GL}(V)$ in $\operatorname{End}(V^{\otimes n})$ spans $B$.
\end{proposition}
\begin{proof}
Denote the span of $g^{\otimes n}$, $g\in\operatorname{GL}(V)$, by $B'$. Let $b\in\operatorname{End}V$ be any element.

We claim that $B'$ contains $b^{\otimes n}$. Indeed, for all values of $t$ but finitely many, $t\cdot\operatorname{Id}+b$ is invertible, so $(t\cdot\operatorname{Id}+b)^{\otimes n}$ belongs to $B'$. This implies that this is true for all $t$, in particular for $t=0$, since $(t\cdot\operatorname{Id}+b)^{\otimes n}$ is a polynomial in $t$.

The rest follows from Lemma 4.56.
\end{proof}
\begin{corollary}[4.59]
As a representation of $S_n\times\operatorname{GL}(V)$, $V^{\otimes n}$ decomposes as $\bigoplus_\lambda V_\lambda\otimes L_\lambda$, where $L_\lambda=\Hom_{S_n}(V_\lambda,V^{\otimes n})$ are distinct irreducible representations of $\operatorname{GL}(V)$ or zero.
\end{corollary}

